\section{Deux positions opposées face aux IAG}\label{sec:positions}

Face à l'irruption des IAG dans les pratiques scolaires, les institutions ont souvent adopté une de deux postures opposées. 
L'une cherche à préserver l'ordre existant en excluant l'outil; l'autre cherche à intégrer l'outil en adaptant les pratiques. 
Ces deux positions méritent d'être présentées dans leur version la plus cohérente, puis soumises à une critique qui révèle leurs limites respectives --- et surtout, la limite qu'elles partagent.

\subsection*{Position 1: l'interdiction}

L'argument central de la position interdictionniste est celui de l'\emph{intégrité académique}. 
Si l'évaluation vise à certifier les compétences d'un individu, toute délégation substantielle de la tâche à un outil est une fraude: non par mauvaise volonté nécessairement, mais parce qu'elle rompt le lien de représentation entre le travail remis et les capacités réelles de son auteur. 
Comme le soulignent les travaux sur l'histoire de l'IA en éducation, les institutions ont toujours cherché à contrôler l'introduction des nouvelles technologies pour préserver la validité de leurs certifications, plutôt que de les remettre en question~\citep{lepage2023recension}. 
L'interdiction s'inscrit dans cette logique de conservation.

Ses arguments sont sérieux. 
La délégation systématique de tâches formatives prive les apprenants du processus même qui construit les compétences: rédiger est formateur non pas parce que cela produit un texte, mais parce que cela oblige à organiser sa pensée, à affronter ses incohérences, à chercher les mots justes. 
Court-circuiter ce processus, c'est obtenir un résultat sans en avoir parcouru le chemin. 
De plus, si les certifications délivrent des diplômes sans correspondance avec des compétences réelles, leur valeur s'érode pour tous leurs titulaires: c'est un problème collectif, pas seulement individuel.

Cependant, la position interdictionniste se heurte à des objections qui ne sont pas seulement pratiques. 
Sur le plan pratique d'abord: les outils de détection de textes générés par IA produisent des taux élevés de faux positifs et ne permettent pas de fonder des décisions académiques fiables~\citep{leroux2015evaluer}. 
L'interdiction risque ainsi de sanctionner inégalement: ceux qui savent mieux dissimuler leur usage sont moins exposés que ceux qui l'utilisent maladroitement. 
Elle creuse donc, paradoxalement, une inégalité supplémentaire là où elle prétend rétablir l'équité.

Mais la limite la plus profonde est épistémologique. 
L'interdiction présuppose que les modalités d'évaluation actuelles --- la dissertation, le mémoire, le rapport écrit --- mesurent effectivement les compétences qu'elles sont censées certifier. 
Or, si ces formes peuvent être satisfaites par une machine, c'est peut-être qu'elles ne mesuraient pas tant une compétence qu'une conformité formelle.
Interdire l'outil sans interroger les critères, c'est préserver l'apparence de la mesure sans en vérifier la validité.

\subsection*{Position 2: l'intégration}

La position intégrationniste repose sur un argument de réalisme pédagogique: les IAG transforment déjà le monde professionnel et les pratiques intellectuelles. 
Ne pas former les apprenants à les utiliser de façon critique, c'est les préparer à un monde qui a déjà changé.%
~\cite{siham2024intelligence} rappellent que chaque grande étape de l'IA en éducation --- des systèmes experts des années 1980 aux systèmes adaptatifs --- a d'abord suscité des craintes avant d'être intégrée, permettant de nouvelles formes de personnalisation et de soutien à l'apprentissage.

Dans sa version la plus réfléchie, cette position ne demande pas seulement d'autoriser les IAG, mais de \emph{repenser les modalités d'évaluation} pour mesurer des compétences moins facilement déléguables: la capacité à évaluer critiquement une production générée par IA, à la corriger, à l'articuler avec un raisonnement propre; l'oral et la mise en situation; le processus de réflexion plutôt que son seul résultat.
Elle anticipe aussi un bénéfice d'équité: les apprenants sans accès à un encadrement familial ou social de qualité pourraient trouver dans les IAG un soutien pour structurer leurs idées, partiellement comparable à ce que d'autres reçoivent de leur environnement culturel.

Cette position a cependant sa propre limite. 
Intégrer les IAG sans redéfinir ce que l'on cherche à évaluer, c'est risquer de certifier des performances qui ne correspondent à aucune compétence réellement acquise. 
La distinction théorique entre IAG comme \emph{substitut} (qui remplace le processus cognitif) et IAG comme \emph{outil} (qui l'amplifie)~\citep{shiohira2021comprendre} est claire en principe, mais indécidable à partir du seul rendu final: l'évaluateur ne voit pas le processus, seulement le produit. 
Une intégration naïve risque donc de produire exactement ce qu'elle prétend éviter: une certification déconnectée des compétences réelles.

\subsection*{Une limite partagée}

Ces deux positions, aussi opposées qu'elles paraissent, partagent un présupposé commun: elles traitent la question des IAG comme un problème de \emph{régulation} de l'accès à un outil, sans interroger la validité des critères d'évaluation sur lesquels elles s'appuient. 
L'une dit \og{}préservons les critères en excluant l'outil\fg{}; l'autre dit \og{}adaptons les outils en conservant l'esprit des critères\fg{}. 
Ni l'une ni l'autre ne demande: ces critères mesurent-ils encore ce qu'ils sont censés mesurer? C'est cette question plus profonde que la section suivante cherche à poser.